\subsection{Вибір та обґрунтування інструментальних засобів розробки}

Реалізація алгоритму двовимірної барицентричної інтерполяції для масштабування зображень вимагає ретельного підбору інструментарію, здатного забезпечити високу швидкість обчислень, ефективну роботу з пам'яттю та гнучкість процесу розробки.

Основною мовою програмування для створення кінцевого програмного продукту обрано C++. Цей вибір фундаментально зумовлений високими вимогами до обчислювальної продуктивності. Обробка растрових матриць із високою роздільною здатністю передбачає виконання мільйонів математичних операцій з плаваючою комою. Мова C++ надає прямий контроль над виділенням та звільненням пам'яті, а також дозволяє виконувати низькорівневу оптимізацію циклів, що є критично важливим для мінімізації часу виконання алгоритму.

Для проектування та реалізації графічного інтерфейсу користувача (GUI) було обрано кросплатформений фреймворк Qt. Цей вибір обґрунтований його повною інтеграцією з мовою програмування C++, високою продуктивністю обробки подій та наявністю вбудованих інструментів для ефективної роботи з графічними даними.

Графічний інтерфейс програми забезпечує зручну взаємодію з користувачем, дозволяючи завантажувати вихідні зображення через діалогові вікна, гнучко задавати коефіцієнти масштабування та візуалізувати отримані результати. Для обробки, конвертації та відображення растрових матриць у реальному часі було використано спеціалізовані класи фреймворку, зокрема \texttt{QImage} та \texttt{QPixmap}. Вони надають оптимізовані механізми для прямого доступу до піксельних даних у пам'яті, що дозволяє мінімізувати накладні витрати під час виведення відмасштабованого зображення на екран.

Архітектурно програмне рішення розділене на два незалежні модулі: обчислювальне ядро та графічний інтерфейс. Взаємодія між ними реалізована через базовий для фреймворку Qt механізм сигналів і слотів. Це дозволило інкапсулювати складні матричні обчислення всередині окремих C++ класів, не перевантажуючи код інтерфейсу користувача. Процес підключення бібліотек Qt та компіляції фінального додатка безшовно інтегрований у єдину конфігурацію CMake.

Для управління процесом збірки проєкту та автоматизації компіляції використано кросплатформену систему CMake. Завдяки CMake реалізовано чітку структуру проєкту з коректним налаштуванням областей видимості змінних та логічним розділенням на підкаталоги. Це значно спрощує інтеграцію сторонніх бібліотек, необхідних для зчитування і збереження графічних файлів, та гарантує легке розгортання проєкту на інших обчислювальних машинах. Безпосередня трансляція вихідного коду у виконуваний файл виконувалася за допомогою стандартизованого компілятора GCC.

Середовище MATLAB застосовувалося виключно на підготовчому етапі для генерації еталонного тестового зображення. Завдяки вбудованим інструментам для роботи з матричними структурами, у MATLAB було сформовано вихідний графічний масив із заданими просторовими характеристиками. Це тестове зображення надалі використовувалося як вхідні дані для верифікації обчислювальної коректності розробленого на C++ алгоритму та аналізу якості фінального масштабування.